\subsection*{分级提示}
\begingroup\small
\begin{enumerate}[nosep,leftmargin=*]
  \item \textbf{口述概念：}seed 管可重放性，标准误管条件于模型的抽样波动，独立基准管实现正确性；bootstrap 的随机性还条件于已经观察到的样本。
  \item \textbf{笔试推导：}均值方差先写协方差双重和；控制变量方差是关于 $\beta$ 的二次函数；bootstrap 小样本可对和的频数做完全枚举。
  \item \textbf{数值编程：}先写出本章公式中的目标量，再显示中间计算。改变参数前预测结果方向，运行后说明差异来自模型、抽样还是数值近似。
  \item \textbf{研究判断：}先识别真正独立的抽样单位，再选 block、cluster 或完整管线重抽样；重要性权重要查支撑、尾部和有效样本量。
\end{enumerate}
\endgroup
